\section*{Chapter \ref{ch:g5:bits-and-bytes} --- Bits and Bytes: Counting in Binary}

\begin{solution}{exo:g5:bits-and-bytes:1}
$5 = 4 + 1$: 00101. $9 = 8 + 1$: 01001. $20 = 16 + 4$: 10100.
\end{solution}

\begin{solution}{exo:g5:bits-and-bytes:2}
$16 + 4 + 2 = 22$.
\end{solution}

\begin{solution}{exo:g5:bits-and-bytes:3}
Four cards: $8 + 4 + 2 + 1 = 15$. Five cards: $16 + 15 = 31$. Eight cards:
$128 + 64 + 32 + 16 + 8 + 4 + 2 + 1 = 255$.
\end{solution}

\begin{solution}{exo:g5:bits-and-bytes:4}
1 \omterm{def:g5:bits-and-bytes:bit}{bit}: 2; 2 \omterm{def:g5:bits-and-bytes:bit}{bits}: 4; 3 \omterm{def:g5:bits-and-bytes:bit}{bits}: 8; 4 \omterm{def:g5:bits-and-bytes:bit}{bits}: 16 patterns.
\end{solution}

\begin{solution}{exo:g5:bits-and-bytes:5}
000, 001, 010, 011, 100, 101, 110, 111.
\end{solution}

\begin{solution}{exo:g5:bits-and-bytes:6}
Five cards reach only 31, which is less than 40, so a sixth card, worth 32, is
needed: $40 = 32 + 8$, written 101000. Six \omterm{def:g5:bits-and-bytes:bit}{bits}.
\end{solution}

\begin{solution}{exo:g5:bits-and-bytes:7}
The \omterm{def:g4:memory-and-storage:byte}{byte} is 10010011: $128 + 16 + 2 + 1 = 147$.
\end{solution}

\begin{solution}{exo:g5:bits-and-bytes:8}
For example: a light switch (on or off); a tiny area of a \omterm{def:g4:memory-and-storage:hard-disk}{hard disk} (magnetised
one way or the other); a cell of \omterm{def:g4:memory-and-storage:flash}{flash memory} (charged or not); a door (open or
shut); a coin on a table (heads or tails).
\end{solution}

\begin{solution}{exo:g5:bits-and-bytes:9}
$0111$ is $4 + 2 + 1 = 7$. Adding 1 gives 8: the cards 4, 2 and 1 turn face
down and the card 8 turns face up. The new number is 1000 in \omterm{def:g5:bits-and-bytes:binary}{binary}, 8 in
ordinary numbers. (It is like $999 + 1 = 1000$: every full place goes back to
0 and carries 1 to the left.)
\end{solution}

\begin{solution}{exo:g5:bits-and-bytes:10}
6 \omterm{def:g5:bits-and-bytes:bit}{bits} make 2, 4, 8, 16, 32, 64 patterns: 64 codes, more than 50. Yes, there
are enough (and 14 to spare).
\end{solution}

\begin{solution}{exo:g5:bits-and-bytes:11}
Each card must be worth one more than all the cards on its right together:
$1 + 2 + 4 = 7$, and the next card is worth 8. Then every number can be shown,
and in one way only. With cards worth 1, 2, 3, 4 and 5, some numbers can be
shown in several ways (3 is the card 3, or $1 + 2$), so a pattern of cards would
not tell a number in a single way; and with all five cards face up, the
largest number would be 15 instead of 31.
\end{solution}
